
The HMM-driven unit selection and WaveNet TTS systems were built from speech at 16 kHz sampling.
Although LSTM-RNNs were trained from speech at 22.05 kHz sampling, 
speech at 16 kHz sampling was synthesized at runtime using a resampling functionality in the Vocaine vocoder \citep{Vocaine}.
Both the LSTM-RNN-based statistical parametric and HMM-driven unit selection speech synthesizers were built from the speech datasets in the 16-bit linear PCM, whereas the WaveNet-based ones were trained from the same speech datasets in the 8-bit $\mu$-law encoding.

The linguistic features include phone, syllable, word, phrase, and utterance-level features \citep{HTS_label} (\eg phone identities, syllable stress, the number of syllables in a word, and position of the current syllable
in a phrase) with additional frame position and phone duration features \citep{Zen_DNN_ICASSP}. These features were derived and associated with speech every 5 milliseconds by phone-level forced alignment at the training stage.
We used LSTM-RNN-based phone duration and autoregressive CNN-based $\log F_0$ prediction models.
They were trained so as to minimize the mean squared errors (MSE).
It is important to note that no post-processing was applied to the audio signals generated from the WaveNets.

The subjective listening tests were blind  and crowdsourced. 
100 sentences not included in the training data were used for evaluation.
Each subject could evaluate up to 8 and 63 stimuli for North American English and Mandarin Chinese, respectively.
Test stimuli were randomly chosen and presented for each subject.
In the paired comparison test, each pair of speech samples was the same text synthesized by the different models.
In the MOS test, each stimulus was presented to subjects in isolation.
Each pair was evaluated by eight subjects in the paired comparison test, and each stimulus was evaluated by eight subjects in the MOS test.
The subjects were paid and native speakers performing the task.
Those ratings (about 40\%) where headphones were not used were excluded when computing the preference and mean opinion scores.
\tblref{tab:sxs} shows the full details of the paired comparison test shown in \figref{fig:sxs2}.

\begin{table}[htbp]
    \centering
    \begin{tabularx}{0.9\textwidth}{L|RRRRR|R}
    \toprule
                       & \multicolumn{5}{|c|}{\textbf{Subjective preference (\%) in naturalness}} &  \\ \cmidrule{2-6}
                       & & & \textbf{WaveNet} & \textbf{WaveNet} & No & \\
     \textbf{Language} & \textbf{LSTM} & \textbf{Concat} & (L) & (L+F) & preference & $p$ value \\ \midrule\midrule
     North    & 23.3 & \textbf{63.6} & & & 13.1 & $\ll 10^{-9}$ \\
     American & 18.7 & & \textbf{69.3} & & 12.0 & $\ll 10^{-9}$ \\
     English  & 7.6 & & & \textbf{82.0} & 10.4 & $\ll 10^{-9}$ \\
              & & 32.4 & \textbf{41.2} & & 26.4 & $0.003$ \\
              & & 20.1 & & \textbf{49.3} & 30.6 & $\ll 10^{-9}$ \\
              & & & 17.8 & \textbf{37.9} & 44.3 & $\ll 10^{-9}$ \\
     \midrule
     Mandarin & \textbf{50.6} & 15.6 & & & 33.8 & $\ll 10^{-9}$ \\
     Chinese  & 25.0 & & 23.3 & & 51.8 & 0.476 \\
              & 12.5 & & & \textbf{29.3} & 58.2 & $\ll 10^{-9}$ \\
              & & 17.6 & \textbf{43.1} & & 39.3 & $\ll 10^{-9}$ \\
              & & 7.6 & & \textbf{55.9} & 36.5 &  $\ll 10^{-9}$ \\
              & & & 10.0 & \textbf{25.5} & 64.5 &  $\ll 10^{-9}$ \\
     \bottomrule
   \end{tabularx}%
   \caption{Subjective preference scores of speech samples between LSTM-RNN-based statistical parametric (\textbf{LSTM}), HMM-driven unit selection concatenative (\textbf{Concat}), and proposed WaveNet-based speech synthesizers.
   Each row of the table denotes scores of a paired comparison test between two synthesizers.
   Scores of the synthesizers which were significantly better than their competing ones at $p < 0.01$ level were shown in the bold type. 
   Note that \textbf{WaveNet} (L) and \textbf{WaveNet} (L+F) correspond to WaveNet conditioned on linguistic features only and that conditioned on both linguistic features and $F_0$ values. }
   \label{tab:sxs}
\end{table}%
